\documentclass[11pt]{article}

% Replace this generic template with the target conference/journal template when needed.
\usepackage[margin=1in]{geometry}
\usepackage{amsmath,amssymb}
\usepackage{booktabs}
\usepackage{graphicx}
\usepackage{multirow}
\usepackage{xcolor}
\usepackage{hyperref}

\newcommand{\method}{SpecEmbedding-Rerank}
\newcommand{\zs}{\mathbf{z}_s}
\newcommand{\zm}{\mathbf{z}_m}
\newcommand{\candidates}{\mathcal{C}}

\title{
  Candidate-Constrained Molecular Reranking for MS/MS-Based Metabolite Annotation
}

\author{
  Anonymous Authors
}

\date{}

\begin{document}

\maketitle

\begin{abstract}
Tandem mass spectrometry (MS/MS) is a central technology for metabolomics, but
accurately annotating unknown spectra remains challenging due to the structural
ambiguity of molecular candidates. We study MS/MS-based molecular annotation as
a bioinformatics data mining problem, where the goal is to retrieve and rank
candidate molecules from large chemical spaces. We propose \method{}, a
two-stage framework that first learns aligned spectrum and molecule embeddings
for coarse retrieval, and then applies a lightweight listwise reranker to refine
the top-ranked molecular candidates. The reranker combines spectrum embeddings,
molecule embeddings, base retrieval scores, and rank-aware candidate context to
distinguish hard candidates with similar mass or formula constraints. Experiments
on metabolomics benchmark datasets show that the proposed reranking stage
improves top-$K$ accuracy and mean reciprocal rank over the base retrieval model.
We further analyze candidate recall upper bounds, candidate construction
protocols, and the contribution of each reranking component.
\end{abstract}

\section{Introduction}

Untargeted metabolomics relies on MS/MS spectra to identify small molecules in
complex biological samples. However, many spectra cannot be confidently
annotated because different molecules may share similar precursor masses,
fragmentation patterns, or molecular formulas. This makes metabolite annotation a
large-scale bioinformatics data mining problem: given an observed spectrum, the
method must mine a large chemical search space and prioritize biologically
plausible molecular structures.

Recent deep learning methods learn neural representations for spectra or
molecules and use embedding similarity for retrieval. While such methods are
efficient for large-scale search, the top retrieved candidates often contain
hard negatives that are structurally or chemically close to the true molecule.
This motivates a second-stage ranking model that focuses on fine-grained
discrimination among a small candidate set.

In this work, we propose \method{}, a candidate-constrained molecular reranking
framework for MS/MS-based metabolite annotation. The framework first uses a
spectrum-molecule alignment model to retrieve top-$K$ candidates. It then
applies a listwise reranker that re-scores these candidates using pairwise
spectrum-molecule features, the base retrieval score, the base rank, and
candidate-candidate context.

\paragraph{Contributions.}
Our main contributions are:

\begin{itemize}
  \item We formulate MS/MS molecular annotation as a candidate-constrained
  bioinformatics data mining problem with explicit candidate recall upper bounds.
  \item We introduce a spectrum-molecule representation learning framework for
  coarse molecular retrieval from MS/MS spectra.
  \item We design a listwise molecular reranker that refines hard top-$K$
  candidates without generating new molecules or accessing the full molecular
  database during reranking.
  \item We provide ablation studies on reranker architecture, feature design,
  loss functions, and candidate construction protocols.
\end{itemize}

\section{Related Work}

\paragraph{MS/MS-Based Molecular Annotation.}
Discuss traditional library search, fragmentation-tree-based methods, molecular
fingerprint prediction, and recent neural approaches for metabolite annotation.
Emphasize the challenge of structural ambiguity and hard candidates.

\paragraph{Representation Learning for Spectra and Molecules.}
Review methods that learn spectrum embeddings, molecular graph embeddings, or
joint spectrum-molecule representations. Position the base model as a retrieval
component that supports scalable candidate search.

\paragraph{Learning to Rank in Bioinformatics Data Mining.}
Discuss pointwise, pairwise, and listwise ranking methods. Explain why molecular
annotation benefits from listwise reranking: candidates for the same spectrum are
not independent, and their relative order carries useful information.

\section{Problem Formulation}

Let $s$ denote an MS/MS spectrum and $m$ denote a candidate molecule represented
by its molecular graph or SMILES string. Given a candidate molecular database
$\mathcal{D}$, the goal is to rank molecules such that the true molecule
$m^{+}$ is placed as high as possible:

\begin{equation}
  \operatorname{rank}(m^{+} \mid s, \mathcal{D}) \rightarrow 1.
\end{equation}

In practice, molecular annotation is often performed under candidate constraints
derived from experimental metadata, such as precursor mass or molecular formula.
We therefore define a candidate set $\candidates_s \subseteq \mathcal{D}$ for
each spectrum $s$. Our reranking task is to improve the order of candidates
inside $\candidates_s$, rather than recover molecules missing from the candidate
set. This distinction is important because the final performance is upper-bounded
by the recall of the initial candidate generation stage.

\section{Method}

\subsection{Overview}

\method{} follows a two-stage retrieval and reranking pipeline. First, a base
spectrum-molecule alignment model maps spectra and molecules into a shared
embedding space and retrieves top-$K$ candidates by similarity. Second, a
molecular reranker refines the order of these candidates using listwise candidate
context.

\begin{figure}[t]
  \centering
  \fbox{
    \parbox{0.88\linewidth}{
      \centering
      MS/MS spectrum $\rightarrow$ spectrum-molecule alignment model
      $\rightarrow$ top-$K$ candidates
      $\rightarrow$ listwise molecule reranker
      $\rightarrow$ reranked candidate list
    }
  }
  \caption{
    Overview of \method{}. The reranker only reorders the initial top-$K$
    candidates and does not generate new molecular structures.
  }
  \label{fig:overview}
\end{figure}

\subsection{Spectrum-Molecule Alignment Model}

The base model encodes an MS/MS spectrum $s$ into a vector $\zs$ and a molecule
$m$ into a vector $\zm$. The retrieval score is computed by normalized embedding
similarity:

\begin{equation}
  b(s, m) = \zs^{\top}\zm.
\end{equation}

The top-$K$ candidates with the highest base scores are retained for reranking.
This stage provides both the candidate molecules and their base scores and ranks.

\subsection{Candidate Reranker}

For each candidate molecule $m_i$ in the top-$K$ set, we construct a pairwise
feature vector:

\begin{equation}
  \mathbf{x}_i =
  [\zs,\ \mathbf{z}_{m_i},\ \zs \odot \mathbf{z}_{m_i},\
  |\zs - \mathbf{z}_{m_i}|,\ b_i,\ \mathbf{e}(r_i)],
\end{equation}

where $b_i$ is the base retrieval score, $r_i$ is the base rank, and
$\mathbf{e}(r_i)$ is a learnable rank embedding. The feature vector is projected
by a multilayer perceptron:

\begin{equation}
  \mathbf{h}_i = \operatorname{MLP}(\mathbf{x}_i).
\end{equation}

To model interactions among candidates for the same spectrum, the sequence
$[\mathbf{h}_1,\ldots,\mathbf{h}_K]$ is passed into a Transformer encoder:

\begin{equation}
  [\tilde{\mathbf{h}}_1,\ldots,\tilde{\mathbf{h}}_K]
  = \operatorname{Transformer}([\mathbf{h}_1,\ldots,\mathbf{h}_K]).
\end{equation}

The reranker predicts a residual score $\Delta_i$ for each candidate and combines
it with the base retrieval score:

\begin{equation}
  s_i = \alpha b_i + \Delta_i,
\end{equation}

where $\alpha$ is a learnable scalar. The final candidate order is obtained by
sorting $\{s_i\}_{i=1}^{K}$ in descending order.

\subsection{Training Objective}

Given the index $y$ of the true molecule within the candidate set, we optimize a
listwise cross-entropy loss:

\begin{equation}
  \mathcal{L}_{\text{CE}}
  = -\log \frac{\exp(s_y)}{\sum_{i=1}^{K}\exp(s_i)}.
\end{equation}

We optionally add a pairwise ranking loss that encourages the positive molecule
to score higher than negative candidates:

\begin{equation}
  \mathcal{L}_{\text{pair}}
  = \frac{1}{K-1}\sum_{i \neq y}
  \log \left(1 + \exp(s_i - s_y + \gamma)\right),
\end{equation}

where $\gamma$ is a margin. The final objective is:

\begin{equation}
  \mathcal{L} = \mathcal{L}_{\text{CE}} +
  \lambda \mathcal{L}_{\text{pair}}.
\end{equation}

\section{Experiments}

\subsection{Datasets}

Describe the metabolomics datasets used in this work, such as MassSpecGym, GNPS,
MoNA, MassBank, or NPLIB1. Report the number of spectra, unique molecules,
training/validation/test splits, and candidate set statistics.

\begin{table}[t]
  \centering
  \caption{Dataset statistics.}
  \label{tab:dataset}
  \begin{tabular}{lrrrr}
    \toprule
    Dataset & Train spectra & Val spectra & Test spectra & Unique molecules \\
    \midrule
    MassSpecGym & -- & -- & -- & -- \\
    GNPS & -- & -- & -- & -- \\
    MoNA & -- & -- & -- & -- \\
    \bottomrule
  \end{tabular}
\end{table}

\subsection{Candidate Construction Protocol}

Clearly specify how candidate sets are constructed. Distinguish official
benchmark candidates, mass-constrained candidates, formula-constrained
candidates, and external database candidates. For fair evaluation, candidate
sets used for validation and testing should not be generated from test labels
unless the experiment is explicitly described as an oracle candidate setting.

\subsection{Baselines}

Compare against representative baselines:

\begin{itemize}
  \item Traditional spectral similarity or library search baselines.
  \item Neural spectrum embedding baselines.
  \item Spectrum-molecule retrieval without reranking.
  \item Pointwise reranking without candidate-candidate context.
\end{itemize}

\subsection{Evaluation Metrics}

We evaluate molecular annotation performance using Top-$K$ accuracy, mean
reciprocal rank (MRR), and candidate recall upper bound. When available, we also
report structure-aware metrics such as MCES for the top-ranked prediction.

\subsection{Implementation Details}

Report encoder dimensions, batch size, optimizer, learning rate, number of
epochs, candidate top-$K$, reranker hidden dimension, Transformer layers, and
hardware. Include random seeds and early stopping criteria for reproducibility.

\subsection{Main Results}

\begin{table}[t]
  \centering
  \caption{Main retrieval and reranking results.}
  \label{tab:main_results}
  \begin{tabular}{llrrrr}
    \toprule
    Dataset & Method & Top-1 & Top-5 & Top-10 & MRR \\
    \midrule
    MassSpecGym & Base retrieval & -- & -- & -- & -- \\
    MassSpecGym & \method{} & -- & -- & -- & -- \\
    \bottomrule
  \end{tabular}
\end{table}

\section{Ablation Studies}

\subsection{Effect of Reranking}

Compare the base retrieval model with the full reranking pipeline to quantify
the gain from the second-stage model.

\subsection{Reranker Architecture}

Compare the Transformer reranker with a pointwise reranker. Vary the number of
Transformer layers and attention heads to analyze the value of candidate
context.

\subsection{Feature Ablation}

Remove individual input features from the reranker, including base score, base
rank embedding, element-wise product, and absolute difference. This identifies
which signals contribute most to hard candidate discrimination.

\subsection{Loss Function Ablation}

Compare listwise cross-entropy alone with the combined cross-entropy and pairwise
ranking objective.

\subsection{Candidate Set Analysis}

Analyze the effect of candidate type and size, including mass candidates,
formula candidates, official benchmark candidates, and different values of
top-$K$. Report the initial candidate recall upper bound to separate retrieval
failure from reranking failure.

\begin{table}[t]
  \centering
  \caption{Candidate set analysis.}
  \label{tab:candidate_analysis}
  \begin{tabular}{llrrrr}
    \toprule
    Candidate type & Top-$K$ & Upper bound & Base Top-1 & Rerank Top-1 & MRR \\
    \midrule
    Mass & 128 & -- & -- & -- & -- \\
    Mass & 256 & -- & -- & -- & -- \\
    Formula & 128 & -- & -- & -- & -- \\
    Formula & 256 & -- & -- & -- & -- \\
    \bottomrule
  \end{tabular}
\end{table}

\section{Case Study}

Provide qualitative examples where the reranker improves the rank of the true
molecule. Discuss whether the corrected candidates share similar mass, formula,
or substructures with the true molecule. This section helps connect the ranking
improvement to chemically meaningful distinctions.

\section{Discussion}

Discuss the strengths and limitations of \method{}. In particular, emphasize
that reranking improves the order of retrieved candidates but cannot recover the
true molecule if it is absent from the initial candidate set. Also discuss the
importance of non-oracle candidate construction protocols for fair evaluation in
bioinformatics data mining.

\section{Conclusion}

We presented \method{}, a two-stage framework for MS/MS-based molecular
annotation. By combining spectrum-molecule embedding retrieval with a listwise
molecular reranker, the method improves hard candidate ranking under mass or
formula constraints. Future work will focus on stricter non-oracle candidate
generation, cross-dataset generalization, and integration with larger molecular
databases.

\bibliographystyle{plain}
% Add references to paper/references.bib and uncomment the next line if needed.
% \bibliography{references}

\end{document}
